\begin{table*}[t]
\centering\footnotesize
\setlength{\tabcolsep}{4pt}
\caption{동일 YOLO의 구성 요소 절제. 국소 보정에 치수·회전 대응을 추가한 효과를 분리한다.}\label{tab:ablation_results}
\begin{tabularx}{\textwidth}{Yrrrrrrr}
\toprule
구성 & \shortstack{코너 중앙값\\(px)} & \shortstack{코너 P90\\(px)} & \shortstack{\pckten\\(\%)} & \shortstack{위치 중앙값\\(cm)} & \shortstack{위치 P90\\(cm)} & \shortstack{회전 중앙값\\(도)} & \shortstack{회전 P90\\(도)} \\
\midrule
N0: 영상 보정 & 5.944 & 42.872 & 67.534 & 7.260 & 38.391 & 2.176 & 85.973 \\
N1: 영상+회전 대응 & 5.921 & 42.513 & 67.560 & 7.274 & 39.128 & 2.183 & 85.989 \\
N2: 영상+치수 & 5.778 & 42.459 & 68.587 & 7.011 & 38.108 & 2.090 & 85.819 \\
N3: 제안 보정 & 5.778 & 42.134 & 68.587 & 7.068 & 37.809 & 2.070 & 85.919 \\
\bottomrule
\end{tabularx}
\tabnote{319장·8코너 평가, 세 seed의 통계 평균. 각 구성의 자세 산출은 319/319이다. N3−N2의 변화는 지표별로 다르며, 치수 정보와 추가 파라미터 효과도 완전히 분리된 것은 아니다. 세부 방향각·겹침·정규화 영상 오차는 보충자료에 보존한다.}
\end{table*}
